% ============================================================
% 最小可编译 LaTeX 论文模板（贴合国赛真实优秀论文惯例，v2.1）
% 用途：避免编译失败重试浪费 token，从此模板开始撰写
% 编译：xelatex paper.tex （两次）
%
% ★ 导言区经完整论文（88 页）实战编译验证（2026-08 批次 2026-08-B1）：
%   - 代码排版用 fvextra + \VerbatimInput 直引源文件（中文注释 100% 保留）；
%     不要用 listings（丢中文注释）或 minted（依赖 Pygments，易崩）
%   - 特殊符号（★△✓∈）经 xeCJK 子分区路由到 Segoe UI Symbol，零缺字形
%   - 行距 1.4（LaTeX 默认偏紧，是"一眼不像国赛论文"的主因之一）
%   - 页码页脚居中、全文无页眉（pagestyle=plain）
%
% 结构依据：组委会官方展示优秀论文的主流结构：
%   摘要页(标题+摘要+关键词，无目录、无队名日期)
%   → 问题重述 → 问题分析 → 模型假设 → 符号说明
%   → 模型的建立与求解（按问题组织，推导内联正文）
%   → 灵敏度与误差分析 → 模型评价 → 参考文献 → 附录
% ============================================================

\documentclass[12pt,a4paper]{ctexart}

% ===== 页面布局（上下 2.54cm、左右 3.18cm，官方页边距 ≥2.5cm） =====
\usepackage[top=2.54cm,bottom=2.54cm,left=3.18cm,right=3.18cm]{geometry}
\usepackage{setspace}
\linespread{1.4}

% ===== 字体（实战验证配置） =====
% 西文/数字用 Times 系（对齐国赛标杆论文观感；无 Times New Roman 的 Linux
% 环境改 \setmainfont{TeX Gyre Termes}）；数学用 newtxmath 与西文同族——
% 不设两者时西文回落 Latin Modern、数学回落 Computer Modern，偏瘦且与国赛观感不符
\setmainfont{Times New Roman}
% 特殊符号（★△✓∈）归入 CJK 类经中文字体回退；其中 ✓ 各中文字体均缺字形，
% 经 xeCJK 子分区机制单独路由到 Segoe UI Symbol（零缺字形验证）
\xeCJKDeclareCharClass{CJK}{"2605, "25B3, "2713, "2208}
\xeCJKDeclareSubCJKBlock{SYM}{"2713 -> "2713}
\setCJKmainfont{SimSun}[AutoFakeBold=2.5, SYM=Segoe UI Symbol]
\setCJKsansfont{SimHei}
\setCJKmonofont{SimSun}[SYM=Segoe UI Symbol]
\setmonofont{Consolas}[Scale=0.82]  % 附录代码等宽字体；中文注释经 setCJKmonofont 回退

% ===== 标题编号（一级"一、"黑体三号居中，二级"（一）"黑体四号，三级阿拉伯小四） =====
\ctexset{
  section = {
    format = \zihao{3}\heiti\centering,
    name = {,、},
    number = \chinese{section},
    aftername = \quad
  },
  subsection = {
    format = \zihao{4}\heiti,
    name = {（,）},
    number = \chinese{subsection},
    aftername = \quad
  },
  subsubsection = {
    format = \zihao{-4}\heiti,
    number = \arabic{subsubsection},
    aftername = \quad
  }
}

% ===== 宏包 =====
\usepackage{amsmath}
\usepackage{newtxmath}  % Times 系数学字体（自带 AMS 符号，禁止再加载 amssymb——会冲突）
\usepackage{booktabs}
\usepackage{graphicx}
\usepackage{float}
\usepackage{caption}
\usepackage{enumitem}
\usepackage{multirow}
\usepackage{fancyhdr}
\usepackage{fvextra}
\usepackage{verbatim}

% ===== 页脚居中页码、无页眉（官方：页码页脚中部，摘要页起编 1） =====
\pagestyle{plain}
\renewcommand{\headrulewidth}{0pt}

% ===== 图表标题五号宋体 =====
\captionsetup{font=small,labelsep=quad}

% ===== 附录代码直引（VerbatimInput 保证与实际代码逐行一致） =====
% 注：breaklines/breakanywhere 属 fvextra 扩展（fancyvrb 原生不支持），必须用 fvextra 承载
\fvset{frame=single,numbers=left,numbersep=6pt,fontsize=\zihao{-5},
       baselinestretch=1.05,tabsize=4,breaklines=true,breakanywhere=true}

% ===== 图表路径 =====
\graphicspath{{figures/}}

\begin{document}

% ============================================================
% 摘要专用页（官方硬约束：≤1 页；只含标题+摘要+关键词）
% 禁止：\maketitle（会带入队名/日期）、\tableofcontents（官方明确不要目录）
% ============================================================
\begin{titlepage}
\thispagestyle{plain}  % 官方：页码从摘要页起编（titlepage 默认 empty 会吞掉页码）
\centering
{\zihao{3}\heiti [论文标题：准确简练醒目]\par}
\vspace{0.5cm}
{\zihao{4}\heiti 摘\hspace{2em}要\par}
\vspace{0.2cm}

\begin{minipage}{0.95\textwidth}
\zihao{-4}
\setlength{\parindent}{2em}
[第 1 段：用 1-3 句交代问题背景与总体思路。]

[以下按问题逐段：针对问题X，建立了什么模型（方法名），用什么工具/算法求解，
得到了什么结论——结论必须定量化（具体数值、百分比、最优方案参数）。
每问一段；若某问含两个子任务，可在段内分述。]

\vspace{0.3cm}
\noindent{\heiti 关键词：}\textbf{[关键词1]}；\textbf{[关键词2]}；\textbf{[关键词3]}；\textbf{[关键词4]}（3-5 个，含核心方法与核心指标；词汇统一加粗，惯例 二.33）
\end{minipage}
% 摘要单页机查锚点（2026-08-B3）：check_latex_compliance.py 读 .aux 中本标签页码，
% 页码=1 则摘要未超页；>1 则溢出（官方硬约束：摘要 ≤1 页）。禁止删除/改名。
\label{abstract:end}
\end{titlepage}

% ============================================================
% 一、问题重述
% ============================================================
\section{问题重述}

\subsection{问题背景}

[问题背景描述。若题目含专有名词/附件说明，可在此处用三线表给出"名词解释与附件说明"。]

\subsection{具体问题重述}

\textbf{问题1}：[用自己的语言提炼为明确的数学任务，避免照搬原文]

\textbf{问题2}：[同上]

\textbf{问题3}：[同上]

% ============================================================
% 二、问题分析
% ============================================================
\section{问题分析}

% 小节标题用"研究内容导向"（如"（一）XX 指标与 YY 参数的关联机理分析"），
% 不用"问题1分析"这类标签式标题。可在此处放机理图、技术路线图、流程图。

\subsection{[问题1研究内容导向的小节标题]}

[该问的核心难点、拟采用的方法路线及理由、与其他问的关联，200-400 字]

% ============================================================
% 三、模型假设
% ============================================================
\section{模型假设}

% 3-6 条，每条为纯断言句（惯例 二.12）。禁止在假设条目后附"该假设是题目隐含""放宽可引入……"
% 类解释跟句——理由/放宽影响的解释归位问题分析节或灵敏度分析节。
\begin{enumerate}[label=\arabic*., leftmargin=2em]
    \item 假设内容（纯断言句）；
    \item 假设内容；
    \item 假设内容。
\end{enumerate}

% ============================================================
% 四、符号说明
% ============================================================
\section{符号说明}

% 正文统一单表（符号/说明/单位），与术语表.json 一致（惯例 二.34）。
% 可按"决策变量/参数"分组；算法术语等细节类表移入附录（见附录"术语补充说明"节），
% 并在此处指明：完整算法术语表见附录。
% 说明须准确具体（如"单位 X 的 Y 成本"），禁止笼统"价格""系数"。
\begin{table}[H]
\centering
\caption{符号说明}
\begin{tabular}{ccc}
\toprule
符号 & 说明 & 单位 \\
\midrule
$x_1$ & [准确具体的含义] & [单位，无量纲则填 /] \\
$x_2$ & [含义] & [单位] \\
\bottomrule
\end{tabular}
\end{table}

% ============================================================
% 五、模型的建立与求解
% ============================================================
% 主体按问题组织：每个问题内部依次完成 建立→求解→验证→结果分析。
% 推导（公式由来、参数定义、模型性质）必须内联在正文对应位置，
% 不得后置到附录；附录仅放代码、支撑材料列表与确实过长的补充推导。
\section{模型的建立与求解}

\subsection{[问题1研究内容导向的小节标题]}

\subsubsection{数据探索与初步判断}
% 条件触发：数据驱动/统计类问题必写（散点图、分布、异常点识别与处理）；
% 纯优化类问题可省略或改为"问题结构分析"。

[数据观察 → 初步规律 → 异常点处理，附图表]

\subsubsection{模型建立}

[模型假设如何落到数学形式。公式在此内联推导：先给形式，再逐项解释参数含义。]
\begin{equation}
y = f(x_1, x_2, \ldots; \theta)
\end{equation}
其中 $\theta$ 为[参数含义]；$x_i$ 为[含义]。

[理论性质（条件触发）：凸性/复杂度/最优性条件（优化类）；显著性检验/拟合优度（统计类）]

\subsubsection{模型求解}

[算法选择：≥2 种候选对比 + 选择理由（≥1 段）。]

% 条件触发：≥2 个模型/算法候选对比时，用下表收敛对比结果（优化口径）
\begin{table}[H]
\centering
\caption{模型/算法对比（优化口径）}
\label{tab:model_compare}
\begin{tabular}{lcccc}
\toprule
模型/算法 & 目标值 & 最优性指标 & 求解耗时 & 备注 \\
\midrule
{MIP 精确解} & {数值} & MIPGap = {x}\% & {秒} & {小规模/精确} \\
{启发式算法} & {数值} & 与精确解偏差 {x}\% & {秒} & {大规模/实时} \\
\bottomrule
\end{tabular}
\end{table}

[求解过程与结果。]

\subsubsection{结果验证与结果分析}

[最优性验证（优化类：MIPGap/枚举完备性/启发式多次运行统计；
统计类：残差分析/拟合优度/显著性检验）。]

% 条件触发：统计拟合/回归预测类问题（含随机性数据的模型），在结果验证处展开以下
% "模型检验"骨架；优化类用最优性验证（MIPGap/枚举完备性），不适用本块。
% 分层说明：残差分析/拟合优度/显著性检验为评审红线最低必查；交叉验证/Bootstrap
% 为可选增强手段（非红线）。以下为写法骨架，非强制全列：
%   a) 残差分析：残差 vs 拟合值散点图，判断同方差性与无系统模式
%   b) 拟合优度：R^2 与调整 R^2（多变量时用后者，避免单纯靠变量数抬升）
%   c) 置信区间：关键参数的置信区间（与显著性检验配合报告）
%   d) 交叉验证：K 折交叉验证（K=5 或 10），报告平均预测误差
%   e) Bootstrap：重采样估计参数/预测区间（次数 ≥1000，报告区间而非点值）

[模型验证（条件触发：含随机性/不确定性/稳态假设时，蒙特卡洛仿真：
仿真次数、收敛情况、与解析结果偏差、偏差 >5\% 的原因分析）。
最优性验证与模型验证是两个独立维度，不可相互替代。]

[结果表格 + 图表引用 + ≥150 字解读（趋势/峰值/拐点 + 模型对应关系 + 实际意义）]

% ============================================================
% 六、灵敏度与误差分析
% ============================================================
\section{灵敏度与误差分析}

\subsection{灵敏度分析}

[核心参数 ±5\%/±10\%/±20\% 扰动，Tornado 图 + 参数敏感度排序。
扰动范围必须与代码实现一致。]
% 图幅宽度统一 0.9\textwidth（多面板图可至 0.95；<0.8 触发机查 WARN，见 LaTeX格式规范 §四）
\begin{figure}[H]
\centering
\includegraphics[width=0.9\textwidth]{fig_sensitivity.pdf}
\caption{参数敏感度 Tornado 图}
\label{fig:tornado}
\end{figure}

% 条件触发：含灵敏度分析时，下表与 Tornado 图互补（表精读 + 图全览）
\begin{table}[H]
\centering
\caption{参数灵敏度汇总表}
\label{tab:sensitivity}
\begin{tabular}{ccccc}
\toprule
参数 & 基准值 & 变化幅度 & 结果变化率 & 灵敏度排序 \\
\midrule
{参数1} & {基准值} & ±5\% & +{x}\% & 1（最敏感） \\
{参数2} & {基准值} & ±5\% & -{x}\% & 2 \\
\bottomrule
\end{tabular}
\end{table}

% 条件触发：两相关图需并排对比时（如不同参数/方法下结果对比），用双子图
\begin{figure}[H]
\centering
\begin{minipage}{0.48\textwidth}
\centering
\includegraphics[width=\textwidth]{fig_a.pdf}\\
{\small（a）情形 A}
\end{minipage}
\hspace{0.04\textwidth}
\begin{minipage}{0.48\textwidth}
\centering
\includegraphics[width=\textwidth]{fig_b.pdf}\\
{\small（b）情形 B}
\end{minipage}
\caption{双子图并排示例（子图用（a）（b）标注，总 caption 概括）}
\label{fig:twin}
\end{figure}

% 条件触发小节（按题型保留，不适用则整节删除）：
\subsection{策略转换临界点分析（条件触发：策略/方案随参数切换的问题）}

[关键参数在何值时最优策略转换，策略转换图或表格。
所有临界点数值必须有求解结果或图表数据支撑，禁止凭直觉。]

\subsection{活跃约束标注（条件触发：约束优化类问题）}

[各情形最优解处各约束的紧/松状态（active/slack），讨论放松紧约束的收益。]

\subsection{误差来源与鲁棒性}

[误差来源 + 量化 + 影响 + 可接受性判断；鲁棒性验证。]

% ============================================================
% 七、模型评价与推广
% ============================================================
\section{模型评价与推广}

\subsection{优势}

[模型优势]

\subsection{局限性}

[模型局限性]

\subsection{改进方向与推广}

[改进方向；可推广到哪些场景]

% 禁令（惯例 一.7）：本节禁止出现"题目设计合理性讨论"独立节与"建议在题目中……"
% 式评点出题人表述。建模阶段的批判性思考（约束完整性/参数合理性）须转化为
% 显式假设与解读策略，落实到模型假设/问题分析节。

% ============================================================
% 参考文献
% ============================================================
\begin{thebibliography}{99}
\bibitem{ref1} 作者. 标题[J]. 期刊, 年, 卷(期): 页码.
\bibitem{ref2} 作者. 标题[M]. 出版社, 年.
% ≥10 篇，GB/T 7714 格式，可查证可访问
\end{thebibliography}

% ============================================================
% 附录（官方：附录页数不限；必须含完整可运行源程序与支撑材料文件列表）
% ============================================================
\appendix

\section{支撑材料文件列表}

% 必须使用三线表（booktabs），两列：文件名 / 文件说明
% 按"代码→结果→其他"分组或按文件类别分组，禁止逐条罗列
\begin{table}[H]
\centering
\caption{支撑材料文件列表}
\begin{tabular}{p{5cm}p{9cm}}
\toprule
\multicolumn{1}{c}{文件名} & \multicolumn{1}{c}{文件说明} \\
\midrule
代码 & [数量] 个 .py 文件（入口是 main.py / problem1\_solver.py，...） \\
结果 & [概述：CSV 结果文件 + 敏感性图 + 可视化图表] \\
{其他文件.pdf} & [文件说明，如"AI 工具使用情况说明"] \\
\bottomrule
\end{tabular}
\end{table}

\section{完整可运行源程序代码}

[建模所用到的全部完整、可运行的源程序代码（官方硬约束：
缺可运行程序或程序运行结果与论文不符，可取消评奖资格）。
用 \VerbatimInput 直接引用 project/ 下源文件，保证与实际代码逐行一致，
禁止"示意性改写"：]

\VerbatimInput{project/problem1\_solver.py}

\section{补充推导（仅限正文确实放不下时）}

[正文模型建立处已给出完整推导链条；此处仅收录过长的中间代数步骤。
若使用本节，正文对应位置必须有明确引用（"详细中间步骤见附录 C"）。
禁止"详细推导较为复杂，从略"类回避性表述。]

\section{术语补充说明}

% 算法术语等细节类表（惯例 二.34）：正文符号说明节只保留统一单表，
% 算法术语表等细节表移入本节，正文符号说明处须指明"完整算法术语表见附录"。
% 若本题无算法术语细节表，可删除本节（并同步删除正文提及）。

\end{document}
